\chapter*{Preface}
\addcontentsline{toc}{chapter}{Preface}

These notes were written while studying machine learning. Each chapter
started life as a self-contained note on a single question --- a
derivation I wanted to understand completely, a concept whose usual
explanation felt incomplete, or a summary of a book chapter. They have
been collected here and organised into parts that follow a natural
progression:

\begin{itemize}
  \item \textbf{Part I} lays down the statistical learning framework:
        risk, empirical risk minimisation, generalisation and evaluation.
  \item \textbf{Part II} gathers the probability and statistics used
        throughout.
  \item \textbf{Part III} covers the linear algebra, geometry and calculus
        behind least squares, SVD, PCA and gradients.
  \item \textbf{Part IV} studies loss functions, likelihood and output units.
  \item \textbf{Part V} treats classical models: perceptron, SVMs, trees,
        boosting and time-series models.
  \item \textbf{Parts VI--VIII} are devoted to neural networks:
        backpropagation, optimisation and regularisation, convolutional
        and recurrent architectures.
  \item \textbf{Part IX} deals with language modelling and structured
        prediction, and \textbf{Part X} with contrastive learning.
  \item \textbf{Part XI} is a longer treatment of machine learning on
        graphs, from harmonic functions and graph Laplacians to graph
        neural networks, kernels and graph embeddings.
\end{itemize}

Because the chapters were written independently they can be read in
almost any order; notation is defined locally in each chapter, and some
topics are deliberately approached from several angles in different
chapters.
